\begin{resumen}
El marco ST-CC descubre patrones de movilidad agrupando regiones y horarios de una matriz de conteos en una rejilla regular, pero cada división se aplica a toda la matriz, de modo que la interacción que ocurre solo dentro de un bloque queda oculta y no existe un criterio estadístico para decidir si un bloque merece subdividirse. Esta tesis propone AST-CC, un árbol que refina la partición de ST-CC mediante pruebas de hipótesis derivadas de la I-divergencia, calibradas con un \emph{bootstrap} de días que respeta la variabilidad entre jornadas, corregidas por la búsqueda del mejor corte y con un reparto del nivel de significancia a lo largo del árbol. Un umbral de tamaño de efecto distingue la estructura relevante de la que solo es detectable por el gran volumen de datos, y entre los cortes admisibles se elige el que logra su ganancia con menos parámetros. En datos sintéticos, el método mantiene la tasa de refinamientos espurios por debajo del 5\,\%, mientras que una nula de Poisson los produce en todas las réplicas, recupera casi por completo la estructura verdadera cuando la partición base agrupa mal y supera a rejillas globales más finas en todas las réplicas con menos parámetros. Sobre los viajes en taxi de Manhattan de noviembre de 2017, el árbol elimina el 95.9 y el 90.7\,\% de la interacción que ST-CC deja dentro de sus bloques en los periodos laborable y feriado, supera a las rejillas globales optimizadas en el periodo laborable y, en el feriado, pierde frente a ellas con la partición de ST-CC y empata cuando parte de una partición mejor. Sus cortes se replican en el resto de 2017 y revelan diferencias de horario entre zonas de volumen parecido, como el desfase del desplazamiento laboral y la separación entre zonas de ocio nocturno y residenciales.
\end{resumen}
